\documentclass[10pt]{article}
\usepackage[margin=20mm]{geometry}
\usepackage{amsmath,amssymb,mathtools,fontspec,unicode-math,xurl,hyperref}
\IfFontExistsTF{DejaVu Serif}{\setmainfont{DejaVu Serif}}{\setmainfont{Latin Modern Roman}}
\setmathfont{latinmodern-math.otf}
\setlength{\parindent}{0pt}\setlength{\parskip}{6pt}\sloppy
\title{Technical review of
A motivic analogue of the \(K(1)\)-local sphere spectrum}
\author{AI technical verification}
\date{}
\begin{document}
\maketitle
Abstract: The assessment of the exact acquired author version is Major. The central conservativity proof uses a false interchange between effective slices and the cofiber of rho. An explicit cellular example with finite slice tower proves that the purported formal identity fails. It does not contradict the conservativity theorem: the additional algebraic-closure vanishing premise is not satisfied by that example. A replacement proof, using module conservativity followed by the orthogonality of a pure slice, restores the required step under the stated hypotheses. A separate extension needed in the topological comparison has not been fully verified and is recorded as a limitation, not as a false equivalence. All twelve physical source pages were read. Scientific adjudication remains pending.\par
\section*{Exact source and theorem scope}
This review concerns M25-0268, A motivic analogue of the \(K(1)\)-local sphere spectrum, by William Balderrama, Kyle Ormsby, and James D. Quigley (the PDF uses J.D. Quigley), DOI `10.4171/jems/1641`. The acquired source is arXiv:2307.13512v3, uploaded on 6 November 2023 at 18:55:39 UTC, with twelve physical pages and matching printed numbering. Its 210,773 bytes have SHA256\par
> a426f34e1ca7d22d734cd8bcf375424a5586827da81696f1a77a1cb88ee496e6.\par
The source manifest has 22,748 bytes and SHA256\par
> 46dad71706ac00cff0f465076b24c159f5708c13f59cf8fc9965be1e625b8e4f.\par
The arXiv record calls v3 the pre-proofs version accepted to JEMS. The EMS landing page reports submission on 10 August 2023, acceptance on 30 October 2023, and online publication on 4 April 2025. The acquired PDF has not been shown to coincide with the publisher text. The frozen title, DOI, and 2025 batch membership remain unchanged.\par
Theorem A treats every base scheme \(S\) on which \(2\) is invertible and every \(X\in\mathrm{SH}(S)\). With \[JQ=\mathop{\mathrm{fib}}\bigl(\psi^3-1:KQ^\wedge_{(2,\eta)}
                  \longrightarrow KQ^\wedge_{(2,\eta)}\bigr),\] it identifies \(L_{KGL/2}X\) with \((X\otimes JQ)^\wedge_{(2,\eta)}\) and the corresponding fiber on completed \(X\otimes KQ\). Here \(KQ\) and \(KGL\) are the absolute motivic spectra specified by the authors; their representability interpretation on regular bases must not be substituted for the definition on all bases. The completion in the result is essential. No uncompleted smashing-localization assertion is inferred.\par
\section*{A false formal step in the conservativity proof}
Equation reference: wrong. Problem 1 (Major: slices do not commute with the rho cofiber as asserted).  In the proof of Theorem 3.1, p. 8, exactness is invoked to identify \[
 s_*(C(\rho)\otimes X/2)\simeq C(\rho)\otimes s_*(X/2).\] A Tate suspension changes the slice index. Consequently this is not a valid consequence of exactness, and the claimed ordinary module identity is false.\par
Equation reference: correct. For any motivic spectrum \(Y\), the correct slice triangle is \[
 \Sigma^{-1,-1}s_{q+1}Y\longrightarrow s_qY
       \longrightarrow s_q(C(\rho)\otimes Y).\] Indeed, \(s_q(\Sigma^{-1,-1}Y)
=\Sigma^{-1,-1}s_{q+1}Y\). Replacing \(s_{q+1}Y\) by \(s_qY\) would incorrectly remove the weight shift. Keeping the slice grading does not solve this problem: the first term requires reindexing, and forgetting that grading does not turn the induced map into ordinary multiplication by \(\rho\) on each slice.\par
\subsection*{A concrete counterexample at the precise formal scope}
Work over \(k=\mathbb Q\), and put \(H={H\mathbb F_2}_{\mathbb Q}\). This is a cellular pure zero-slice. For the cellularity assertion, the acquired Röndigs–Spitzweck–Østvær reference [RSO], arXiv v2, Proposition 3.31 and Lemma 3.32 (physical p. 29), supplies a cell presentation of \(H\mathbb Z_{\mathbb Q}\); taking the cofiber of \(2\) gives \(H\). Its zero-slice property follows from the same reference, Definition 2.1 and Theorem 2.7, physical pp. 5–6.\par
First the triangle \[\Sigma^{-1,-1}H\xrightarrow{\rho}H\longrightarrow C(\rho)H\] gives \[s_0(C(\rho)H)=H,\qquad
 s_{-1}(C(\rho)H)=\Sigma^{0,-1}H,\qquad
 s_q(C(\rho)H)=0\quad(q\ne-1,0).\] Thus even the simple example \(Y=H\) distinguishes the two sides of the asserted interchange after inverting \(\rho\). For an example satisfying more of the structural hypotheses in the surrounding argument, take \[X=C(\rho)H,\qquad Y=X/2.\] Then \(X\) is cellular and \(2\)-torsion, and its slice tower has only the two nonzero slices just computed. It is therefore convergent. Its real-points Betti realization vanishes, since \(\rho\) realizes to an equivalence. These are the relevant structural hypotheses.\par
Equation reference: nonzero. Because \(X\) is an \(H\)-module, \(2\) acts trivially, and hence \[Y\simeq X\oplus\Sigma^{1,0}X,\qquad
 s_1Y=0,\qquad s_0Y=H\oplus\Sigma^{1,0}H.\] Equation [correct] (\#correct) at \(q=0\) now yields \[
 s_0(C(\rho)\otimes Y)=H\oplus\Sigma^{1,0}H.\] The total slice object on the left of [wrong] (\#wrong) therefore has nonzero \(\rho\)-localization.\par
Here is an explicit check of that nonvanishing. In mod-\(2\) Milnor \(K\)-theory, the sign rule \[\{a\}\longmapsto \epsilon(a)t,\qquad
 \epsilon(a)=
 \begin{cases}1,&a<0,\\0,&a>0,\end{cases}\] defines a homomorphism \(k_*^M(\mathbb Q)\to\mathbb F_2[t]\). Multiplicativity of the sign verifies the bilinear relations, and \(\epsilon(a)\epsilon(1-a)=0\) verifies the Steinberg relation. In particular, \(\{-1\}^n\) maps to \(t^n\ne0\) for every \(n\). Using \(\pi_{*,*}H=k_*^M(\mathbb Q)[\tau]\), this proves \(H[\rho^{-1}]\ne0\), as required in [nonzero] (\#nonzero). On the other hand, \[\bigl(C(\rho)\otimes s_*Y\bigr)[\rho^{-1}]=0\] by the definition of \(C(\rho)\). The two ordinary module objects in [wrong] (\#wrong) cannot be equivalent.\par
The scope of this counterexample is crucial. The chosen \(X\) is not killed by pullback to an algebraic closure. It is a counterexample to the purported general formal interchange, including the structural conditions listed above, but it is not a counterexample to the vanishing implication used in Theorem 3.1, to Theorem 3.1 itself, or to Theorem A.\par
\subsection*{Why this affects a core argument}
Theorem 3.1 is the principal conservativity result of Section 3. Its proof uses [wrong] (\#wrong) to treat the slices as \(\rho\)-torsion and apply Proposition 3.8. Equation [nonzero] (\#nonzero) shows directly why this inference is unavailable: slices of a \(\rho\)-torsion object need not themselves be \(\rho\)-torsion. The resulting gap lies in the central reduction to \(\mathbb C\), used by Proposition 3.12 and then Theorem A. The Major assessment concerns this essential proof step; it does not assert a false final theorem or require that the defect be irreparable.\par
\section*{Replacement argument and its boundary}
A pure slice supplies exactly the additional fact needed to repair the argument. If \(M\) is a cellular \({H\mathbb F_2}_k\)-module, is a pure \(q\)-slice, and becomes zero over an algebraic closure of a prime field \(k\ne\mathbb F_2\), then \(C(\rho)\otimes M\) is cellular, \(\rho\)-torsion, and also becomes zero there. Proposition 3.8 gives \(C(\rho)\otimes M=0\). Therefore \[\rho:\Sigma^{-1,-1}M\xrightarrow{\ \simeq\ }M.\] A \(q\)-slice lies in \(\mathrm{SH}(k)^{\mathrm{eff}}(q)\) and is right orthogonal to \(\mathrm{SH}(k)^{\mathrm{eff}}(q+1)\). But the equivalence implies \[M\simeq\Sigma^{1,1}M
      \in\mathrm{SH}(k)^{\mathrm{eff}}(q+1).\] Thus its identity morphism is zero, and \(M=0\).\par
The global order of the argument also matters. It avoids assuming that arbitrary base change preserves an infinite slice-completion limit. Let \(S_0=\operatorname{Spec}\mathbb Z[1/2]\) and let \(X\) satisfy the original hypotheses of Theorem 3.1, with \(X_{\mathbb C}=0\). Put \[Y=C(\rho)\otimes X/2,\qquad N_q=s_qY
 \quad\text{over }S_0.\] By Lemma 3.3, \(X_K=0\) for every algebraic closure \(K\) of a prime field \(k\ne\mathbb F_2\). Cellular slice base change identifies \((N_q)_k\) with \(s_q(Y_k)\) and identifies its pullback to \(K\) with zero. The applicable reference is [RSO], Corollary 2.17, arXiv v2 physical p. 12. Its connected compatible-pair assumptions hold for \(S_0\) with coefficients \(\mathbb Z_{(2)}\), for these prime fields, and for their algebraic closures; see Definition 2.1 and Remark 2.2, physical p. 5.\par
Lemma 3.5 makes \((N_q)_k\) a cellular \({H\mathbb F_2}_k\)-module, since \(Y_k=(C(\rho)\otimes X_k)/2\). Apply the preceding pure-slice argument to obtain \((N_q)_k=0\) for every such \(k\). Lemma 3.2 then gives \(N_q=0\) over \(S_0\). This holds for every \(q\). The original convergence hypothesis on \(X\) passes to \(Y\) through finite cofibers and Tate shifts, so \(Y=0\) over \(S_0\).\par
Finally, for each prime field \(k\ne\mathbb F_2\), \(C(\rho)\otimes(X_k/2)=0\). Lemma 3.4 applies: over finite fields \(\rho\) is nilpotent, and over \(\mathbb Q\) the original real Betti hypothesis supplies the required vanishing. Hence \(X_k/2=0\). Lemma 3.2 gives \(X/2=0\) globally. Since \(X\) is \(2\)-torsion, invertibility of \(2\) on \(X\) forces \(X=0\).\par
This is a complete replacement of the faulty step relative to the named module-conservativity, field comparison, and slice-base-change results. It preserves the original theorem’s hypotheses. It neither treats [wrong] (\#wrong) as correct nor assumes convergence after pullback. No stronger conservativity statement is claimed, and no author adoption of this repair is known.\par
\section*{Other proof checks and an unresolved comparison qualification}
The argument in Section 2 was checked using \(KGL\simeq KQ/\eta\) and the finite \((2,\eta)\) quotient. Completed \(KQ\)-modules supply the stated local objects; testing the unit map after smashing with \(KGL/2\) reduces the arbitrary-\(X\) assertion to the sphere. The use of completion was retained throughout.\par
The base-change claim in Lemma 3.9 concerns the appropriate finite-quotient object: \[KGL/2\otimes JQ\simeq
 KGL/2\otimes\mathop{\mathrm{fib}}\bigl(\psi^3-1:
 KQ[1/3]\longrightarrow KQ[1/3]\bigr).\] It does not require a general assertion that raw completed \(JQ\) commutes with every base change. Real-points Betti realization in Lemma 3.11 was distinguished from the equivariant Betti realization discussed in the introduction. The needed slice-convergence reduction is over \(S_0\); this review does not independently certify the broadest arbitrary-base wording of Lemma 3.10.\par
In Lemma 3.6, \(\tau\)-free is explicitly defined to mean free over \(\mathbb F_2[\tau]\), not merely torsion-free. The iterative construction kills the maps on low homotopy groups, and its colimit supplies the connective approximation. Lemma 3.7 can also be checked directly: construct the stem-zero approximation to \({H\mathbb F_2}_k\) by successively killing negative-stem free cells. A map representing an element of a module concentrated in stem zero extends through each step because the obstruction groups in negative stems vanish. This identifies the approximation with \(E\) and produces the required maps from weight shifts of \(E\). The resulting basis map is an equivalence by cellularity. This check avoids relying on an unstated convergence claim for the displayed universal coefficient spectral sequence. Proposition 3.8 then uses the finite stem range of \(C(\rho)\otimes{H\mathbb F_2}_k\) for the specified prime fields; its prime-field restriction was not omitted.\par
For Section 4, the acquired published GIKR reference [GIKR], physical p. 14 (printed p. 3610), Proposition 3.18, preserves the relevant cofiber sequence when the map on odd \(MU\)-homology is injective. This corresponds to Proposition 3.17 in the cited preprint numbering. The evenness argument was checked: the Wood cofiber sequence and triviality of \(\eta\) on \(MU\) make \(MU\otimes KO\) a retract of \(MU\otimes KU\). The latter has even, torsion-free homotopy, so the mod-\(2\) objects needed here also have even \(MU\)-homology. Thus the stated fiber-sequence use of that result has the required input.\par
There remains a precise limit to the tensor comparison check. GIKR’s Proposition 3.26, physical pp. 16–17 (printed pp. 3612–3613), requires its first variable to be a bounded-below, finite-type even-cell complex. This is Proposition 3.25 in the earlier numbering. The source paper instead applies it to periodic \(KU\), which is not bounded below, after describing \(KU\) as a filtered colimit of finite even-cell complexes. The adjacent GIKR colimit lemma has a uniform bounded-below and increasing-connectivity hypothesis. Consequently the cited statement alone does not literally cover the passage used on source p. 10.\par
No counterexample to \[\Gamma_*(KU)\otimes_{\Gamma_*(\mathbb S)}\Gamma_*(KO)
       \simeq\Gamma_*(KU\otimes KO)\] has been established here. Results for synthetic spectra give a plausible extension route for filtered colimits of finite \(MU\)-projectives, but the full compatibility with this \(\Gamma_*\) construction has not been independently completed. This is recorded as an unverified dependency qualification, not as a second confirmed false statement. A revised proof should give the appropriate extension or a compatible alternative comparison.\par
\section*{Coverage, evidence, and limitations}
All twelve source pages, including the introduction, every lemma and proof, bibliography, and affiliations, were read as text and actually viewed as page images. The review separately checked the exactness argument, the explicit Milnor-\(K\) sign computation, the counterexample’s hypotheses and exclusions, and the replacement pure-slice proof. These are analytic checks; no author’s code was executed and no numerical test is presented as a proof.\par
The primary reviewer read their saved checks, reconciled the counterexample scope, and performed the report assessment. These are AI assistance records, not human reviews or independent acceptance. The runtime self-description and the externally reported UI selection are recorded separately; backend equivalence is not independently verified.\par
The acquired RSO reference pages used above and the three GIKR evidence pages were inspected at their actual formulas. This is not a claim to have independently reproved those entire references. The additional reference checks and their restricted coverage are preserved in the supporting records. Bounded searches of the arXiv and publisher records and for corrections, supplements, and code found no applicable separate correction or computational supplement. This does not establish nonexistence; absence of code is not a defect of this mathematical paper.\par
One intermediate checkpoint prematurely listed source images whose tool response had been truncated. That record was retained, and a later correction records the separate successful viewing of those images. Rendering-time records likewise retain their original false inspection values; actual later viewing is recorded separately.\par
The assessment is therefore Major, at the core-proof-step scope specified in Problem 1, with a complete local replacement argument and the separate comparison limitation above. The acquired version’s main localization theorem has not been disproved. `scientific\_decision` remains `pending`; `human\_reviews` is empty.\par
9 W. Balderrama, K. Ormsby, and J. D. Quigley, A motivic analogue of the \(K(1)\)-local sphere spectrum, arXiv:2307.13512v3, 6 November 2023. https://arxiv.org/abs/2307.13512v3. EMS Press, publication record for DOI 10.4171/jems/1641. https://ems.press/journals/jems/articles/14298698. O. Röndigs, M. Spitzweck, and P. A. Østvær, The first stable homotopy groups of motivic spheres, arXiv:1604.00365v2; subsequently Ann. of Math. (2) 189 (2019), 1–74. The physical page citations above refer to the acquired arXiv v2, not the publisher pagination. https://arxiv.org/abs/1604.00365v2. B. Gheorghe, D. C. Isaksen, A. Krause, and N. Ricka, \(\mathbb C\)-motivic modular forms, J. Eur. Math. Soc. 24 (2022), 3597–3628, DOI 10.4171/JEMS/1171. Acquired published PDF: https://ems.press/content/serial-article-files/32750.\par
\end{document}
